\section{What changed in the five archived proofs}
\label{sec:engineering}

The reference implementation assembles the zero-free region from named
results for zero factorization, a local logarithmic derivative estimate,
the three-height inequality, growth, a reciprocal bound, and bounded
heights. Its entry point is \lean{Formalism.ZF.zeroFreeRegion}.
The challenge assigns its 21-file dependency closure 21,641 normalized
Lean tokens. By the first snapshot studied here, the submitted proof
already used the entire function \(F(s)=(s-1)\zeta(s)\) at all three
heights. The four source comparisons below start from that implementation.

Table~\ref{tab:runs} records the public results. Comparing consecutive
archived team proofs, with v22 omitted, gives savings of 236, 168, 300
and 36 tokens. These differences measure all changes in each patch together.

\subsection{The Dirichlet series at the three centers (v21 to v23)}

In v21 the common quantity \(D(z)\) was defined as
\(-\Re(\zeta'(z)/\zeta(z))\). To prove positivity, the helper
\lean{triple} first converted three convergent Dirichlet series into this
form. In v23, \(D(z)\) is defined directly as the real part of the
von Mangoldt Dirichlet series. The library result
\lean{LSeries_vonMangoldt_eq_deriv_riemannZeta_div}
identifies that series with \(-\zeta'/\zeta\) for \(\Re z>1\).
The conversion now occurs once, in the identity relating \(D\) to
\(F'/F\). At the point of summation,
\lean{LSeriesSummable_vonMangoldt} gives convergence at every required
height. With \(S(v)\) the real \lean{HasSum} statement at \(a+iv\),
the submitted combination reads
\begin{verbatim}
(((S 0).mul_left 3).add ((S s.im).mul_left 4) |>.add (S (2*s.im)))
\end{verbatim}
\lean{HasSum.nonneg} transfers termwise nonnegativity to this sum.
Expanding the cosine of twice the angle turns the trigonometric factor
into \(2(1+\cos\theta)^2\).

The same edit moved the uses of \lean{zeroHeight} and \lean{poleTerm_le}
into the contradiction argument. It also simplified the passage from
the modulus-one character product
\lean{DirichletCharacter.norm_LFunction_product_ge_one} to a product of
values of \(F\). Given the horizontal gap \(d\leq a-1\), the single
inequality \(d\leq |a+iv-1|\) holds for every \(v\). Rewriting
\(\zeta(z)=(z-1)^{-1}F(z)\) in the product exposes all three
denominators to the same order-preserving rewrite. The earlier script
constructed and applied a separate bound on \(|\zeta(a+iv)|\) at each
height. See \texttt{diffs/96267-to-96281.diff}.

\subsection{A derivative bound from one circle (v23 to v24)}

In the zero-free branch of the disk argument, let \(g\) denote the
primitive of \(f'/f\), normalized by \(g(0)=0\), and put
\(A=\log B+1\). Both proofs establish \(\Re g\leq A\) on the disk of
radius \(1/6\). In v23, \lean{HasDerivAt.le_of_lip'} required an
eventual estimate \(|g(w)-g(0)|\leq C|w|\) at zero. The proof therefore
introduced a neighborhood filter and kept the variable \(|w|\) throughout
the estimate. In v24, \lean{Complex.borelCaratheodory_zero} supplies
\[
 |g(w)|\leq \frac{2A|w|}{1/6-|w|}.
\]
Evaluating on \(|w|=1/12\) gives \(|g(w)|\leq 2A\).
The Cauchy estimate
\lean{Complex.norm_deriv_le_of_forall_mem_sphere_norm_le} takes a bound
\(C\) on a sphere of radius \(r\), together with differentiability inside
and continuity on the closure, and returns \(|g'(0)|\leq C/r\).
Here \(C/r=24A\). Once the boundary estimate is in the hypothesis
\lean{H}, the archived script closes its numerical step with
\begin{verbatim}
        simp at hz
        norm_num [hz] at H;linarith
\end{verbatim}

The interface between the disk estimate and the main theorem changed
as well. The helper \lean{good} now returns a complex equality for \(F'/F\)
alongside nonvanishing. The main proof instantiates a single local
function at \(0,t,2t\) and rewrites all three resulting inequalities
together. This removes the separately stated \lean{scaledWindow} and the
real abbreviation \(D\). In the zero-removal branch, monotonicity of
\(\log\) gives \(\log(B/2)\leq\log B\) directly. The source credits
this last observation to Pi's public submission \#96263.
See \texttt{diffs/96281-to-96285.diff}.

\subsection{The induction parameter and the growth bound (v24 to v25)}

Removing a zero \(\rho\) within radius \(1/6\) gives a factorization
\(f(w)=(w-\rho)g(w)\), where the relative growth bound for \(g\) on
the disk of radius \(1/2\) is \(B/2\). In v24 the induction starts with
\(B<2^n\). In v25, \lean{exists_nat_gt} supplies \(n>2B\).
At a successor step \(2B<n+1\) and \(B\geq 1\) imply \(B<n\).
This is precisely the induction premise \(2(B/2)<n\) for the new
function. Linear arithmetic now suffices for the decrease.

The two versions use the same divided difference \lean{dslope}.
The library lemma \lean{differentiableOn_dslope} transfers
differentiability across the removed zero, and
\lean{sub_smul_dslope_of_zero} identifies the product with \(f\)
pointwise. In v25, function extensionality turns the latter into an
equality of functions. Destructuring the existence statement with
\lean{rfl} replaces \(f\) by its product in every hypothesis and in the
goal. In particular, \lean{logDeriv_mul} now applies directly to the
goal. Its hypotheses are nonvanishing and differentiability of both
factors at zero. The maximum-modulus argument still uses
\lean{Complex.norm_le_of_forall_mem_frontier_norm_le} to pass from the
boundary estimate to the closed disk.

Independently, the growth theorem changed its mathematical statement.
The v24 bound was \(|F(z)|\leq R^8\) under \(R\geq2\),
\(\Re z\geq1/2\) and \(|z|\leq4R\). The v25 lemma gives
\(|F(z)|\leq(2|z|+2)^2\) under the sole hypothesis \(\Re z\geq1/2\).
The disk proof specializes it to \((16R)^2\) at all three centers.
The Mellin remainder is now defined by its integral, the form needed
for the norm bound. A function equality identifies it with the Mellin
transform at \(-z\) when differentiability is needed.
\lean{Complex.isOpen_re_gt} supplies the open half-planes used in
analytic continuation. See \texttt{diffs/96285-to-96293.diff}.

\subsection{Two algebraic end steps (v25 to v26)}

In the Abel identity, v26 asks directly for the factored expression
\(\zeta(z)=z((z-1)^{-1}-E(z))\), with \(E(z)\) the remainder integral.
The partial summation formula \lean{LSeries_eq_mul_integral'}
already has \(z\) outside its integral. The new expression therefore
matches the conclusion of integral subtraction without redistributing
this factor.

In the branch of the induction that removes a zero, the recursive inequality
\lean{ih0} contains \(\log(B/2)\) on its right-hand side. From \(1\leq B\),
the proof uses \(B/2\leq B\) inside that inequality, then closes the
linear goal. The two lines in v26 are
\begin{verbatim}
      grw [half_le_self (zero_le_one.trans hB)] at ih0
      linarith
\end{verbatim}
Finally, the summand with \(\theta=t\log n\) and
\(w=\Lambda(n)\exp(-a\log n)\) reduces to
\(w(3+4\cos\theta+\cos(2\theta))\).
Previously a \lean{calc} block proved \(2w(1+\cos\theta)^2\geq0\) and
then proved the polynomial identity by \lean{ring}. In v26,
\lean{linear_combination} multiplies the inequality for the square by
\(2w\) and checks the resulting identity in one step. These changes,
together with the smaller elaboration edits in
\texttt{diffs/96293-to-96295.diff}, account for the final 36-token drop.
